\documentclass[11pt,a4paper]{article}
\usepackage{fontspec,polyglossia}
\setdefaultlanguage{russian}
\setmainfont{Arial}
\newfontfamily\cyrillicfont{Arial}[Script=Cyrillic]
\usepackage[margin=21mm]{geometry}
\usepackage{xcolor,hyperref}
\definecolor{ink}{HTML}{142D43}
\definecolor{teal}{HTML}{007F86}
\hypersetup{colorlinks=true,urlcolor=teal,pdftitle={Заметки докладчика: фемтоскопия},pdfauthor={Балинов Темир}}
\setlength{\parindent}{0pt}
\setlength{\parskip}{6pt}
\linespread{1.08}
\setlength{\emergencystretch}{3em}
\begin{document}
{\LARGE\bfseries\color{ink} От импульсов к пространству}\par
{\large Заметки докладчика}\par
Балинов Темир, физический факультет МГУ. 19.10.2026.\par
28 основных слайдов: 43:45 с объяснением формул и паузами. Резерв: слайды 29–38.\par
\medskip
\section*{Слайд 01. От импульсов к пространству}
\textbf{Время:} 00:30

\textbf{Тезис:} размеры и эволюцию исчезнувшей системы можно исследовать через импульсные корреляции.

\textbf{Визуал:} связать расходящиеся импульсы на титуле с неизвестным источником.

\textbf{Текст:} Представим, что короткое столкновение уже закончилось, а мы хотим узнать размеры возникшей в нём системы. Сегодня выясним, зачем физикам такой вопрос и как пары зарегистрированных частиц помогают на него отвечать. Сначала поймём идею метода, затем посмотрим, как измеряют нужные величины в STAR и как проверяют их интерпретацию с помощью SMASH.

\textbf{Переход:} Почему для этого сталкивают тяжёлые ядра?

\clearpage
\section*{Слайд 02. Зачем сталкивать тяжёлые ядра?}
\textbf{Время:} 02:00

\textbf{Тезис:} столкновение ядер позволяет исследовать расширение плотной системы многих взаимодействующих частиц.

\textbf{Визуал:} начальное столкновение → плотное облако → расширение; связать рисунок с вопросами о размере и движении.

\textbf{Текст:} Тяжёлое ядро содержит много нуклонов — протонов и нейтронов. Сталкивая такие ядра, мы создаём много взаимодействующих частиц в малом объёме. Часть энергии движения превращается в возбуждение и рождение новых частиц. Система расширяется, становится разреженной и распадается на конечные частицы.

Мы хотим понять, насколько она успела увеличиться, как двигалось вещество и как частицы его покидали. Разные модели эволюции могут предсказывать разные размеры области испускания, даже если похожим образом описывают число конечных частиц. Значит, размеры дают дополнительную проверку физической картины.

В нашем рассказе будем использовать пионы — лёгкие адроны, составные частицы сильного взаимодействия. У пиона спин ноль; положительный и отрицательный пионы — разные сорта. Их много, а пары одинакового заряда удобны для исследования корреляций. При высокой энергии возможно образование кварк-глюонной плазмы; для применимости самого метода это не обязательное исходное условие.

Пауза: вернуться к двум вопросам — насколько большая система и как она расширялась?

\textbf{Переход:} Прежде чем измерять размер, уточним, о какой стадии столкновения идёт речь.

\clearpage
\section*{Слайд 03. Какую стадию помнят конечные пионы?}
\textbf{Время:} 01:15

\textbf{Тезис:} фемтоскопия пионов чувствительна к области их последнего взаимодействия и рождения при распадах.

\textbf{Визуал:} ранняя плотная стадия → расширение → разные точки и времена испускания.

\textbf{Текст:} Пион может неоднократно рассеяться и изменить импульс. Поэтому конечный импульс не является простой меткой того места, где он впервые появился.

По мере расширения реакции редеют. Химическое вымораживание связано с прекращением эффективного изменения выходов частиц, кинетическое — с прекращением перерассеяния. Химическое вымораживание не останавливает последующие распады резонансов. Конечный пион мог родиться при таком распаде.

Для разных частиц это происходит в разных точках и в разное время. Источник в нашем дальнейшем расчёте — именно такое распределение испускания, а не мгновенная фотография раннего столкновения. Когда сравниваем радиусы, важно помнить, какую стадию и какой набор частиц мы выбрали.

\textbf{Переход:} Наша область испускания уже исчезла. Какая информация о ней остаётся доступной?

\clearpage
\section*{Слайд 04. Как узнать размеры исчезнувшего источника?}
\textbf{Время:} 01:00

\textbf{Тезис:} наблюдаемые импульсы позволяют поставить обратную задачу о недоступных точках испускания.

\textbf{Визуал:} неизвестная область с точками испускания → известные конечные импульсы.

\textbf{Текст:} Детектор восстанавливает импульсы частиц и помогает определить их сорт. Но мы хотим узнать, откуда и когда частицы покинули источник.

Интересующие расстояния — фемтометры, то есть десять в минус пятнадцатой степени метра. Общая вершина столкновения измеряется, однако её разрешение несопоставимо с разделениями между точками последнего взаимодействия разных пионов. Обычное продолжение треков назад эту задачу не решает.

Значит, нужно искать сведения об источнике в самих конечных импульсах. Не только в каждом по отдельности, но и в том, как часто они встречаются вместе.

\textbf{Переход:} Идею извлекать размер из совместной регистрации удобно сначала увидеть в оптическом опыте.

\clearpage
\section*{Слайд 05. HBT: размер звезды без её фотографии}
\textbf{Время:} 01:45

\textbf{Тезис:} корреляция интенсивностей может нести информацию об угловом размере источника.

\textbf{Визуал:} звезда → два приёмника → сравнение сигналов при разных расстояниях между приёмниками.

\textbf{Текст:} Звезда далека, и её угловой размер мал. Хэнбери Браун и Твисс предложили получать сведения о размере через связь между двумя зарегистрированными сигналами. В 1956 году после лабораторной демонстрации они применили интенсивностную интерферометрию к Сириусу.

Здесь сравниваются флуктуации фототоков двух раздельных светоприёмников. Каждый приёмник получает свет от всей звезды; стрелки на схеме показывают каналы регистрации, а не отдельные участки диска. Меняя расстояние между приёмниками, исследуют, как меняется корреляция интенсивностей. Её пространственный масштаб позволяет оценить угловой размер звезды. Это не обычное изображение диска и не мерцание звезды, которое видно глазом.

Не нужно сейчас осваивать весь оптический расчёт. Запомним конструкцию вопроса: один приёмник сообщает, сколько света пришло; два дают ещё и информацию о совместной регистрации. В этой дополнительной информации обнаруживается размер источника.

Пауза: показать, что меняется в опыте — расстояние между приёмниками, а не сама звезда.

\textbf{Переход:} Как общий принцип парных корреляций проявился у частиц в лаборатории?

\clearpage
\section*{Слайд 06. От фотонов — к пионам, затем к столкновениям ядер}
\textbf{Время:} 01:15

\textbf{Тезис:} общая бозонная симметрия связывает историю оптических и пионных корреляций, но измерительные процедуры различны.

\textbf{Визуал:} три ступени: фотоны и HBT → пионные пары → распределённый источник в ядерном столкновении.

\textbf{Текст:} Пионы, как и фотоны, — бозоны: их спин целый, а квантовое описание симметрично при перестановке одинаковых частиц. В работе Голдхаберов, Ли и Пайса 1960 года рассматривались пионы при аннигиляции антипротона с протоном. Симметризация одинаково заряженных пионов позволяла качественно объяснить наблюдавшиеся угловые корреляции. В модельном расчёте они зависели от размера объёма испускания. Возникала обратная задача: какие размеры совместимы с данными? Это ещё не было точным измерением радиуса тяжёлоионного источника.

Общая основа — свойства тождественных бозонов, то есть частиц с целым спином, и информация в парах. Но телескоп в лабораторию буквально не переносится: в оптике меняли базу приёмников, а здесь сравнивают импульсы двух частиц одного события.

В 1970-е годы Копылов и Подгорецкий развивали связь таких корреляций с пространственными и временными характеристиками испускания. Чтобы воспользоваться этой идеей, сначала научимся отделять совместную регистрацию от простого большого числа частиц.

\textbf{Переход:} Что именно мы будем называть корреляцией и с чем сравнивать число пар?

\clearpage
\section*{Слайд 07. Знать спектр — ещё не знать пары}
\textbf{Время:} 01:30

\textbf{Тезис:} C сравнивает совместную регистрацию с независимым ожиданием; независимому контролю соответствует единица.

\textbf{Визуал:} одночастичный спектр → совместное распределение → смысл C = 1.

\textbf{Текст:} Представим много событий с зарегистрированными частицами. Гистограмма одного импульса показывает, какие значения встречаются часто. Но она не отвечает, какие импульсы встречаются вместе в одном столкновении. Для этого нужно совместное распределение пары.

При независимости оно получается произведением одночастичных распределений с учётом общей нормировки. Корреляционная функция C сравнивает совместную частоту с независимым ожиданием: берём отношение двухчастичной плотности к произведению одночастичных и вводим общую нормировку N. В эксперименте используем плотности выходов частиц и пар; N выбираем так, чтобы независимому фону в контрольной области соответствовало C = 1. Значение выше единицы — выбранные импульсы встречаются вместе чаще контроля; ниже — реже. Нормированный спектр слева — лишь учебная иллюстрация одночастичного распределения.

Можно сохранить отдельные импульсы, но переставить частицы между событиями. Спектр сохранится, а парные связи изменятся. Поэтому пары дают новую информацию. Однако связь ещё не обязательно вызвана размером: существуют распады, законы сохранения и эффекты регистрации.

Пауза: прежде чем искать физический сигнал, нужно понимать значение базовой линии C = 1.

\textbf{Переход:} В какой переменной искать избыток похожих импульсов?

\clearpage
\section*{Слайд 08. Какие пары считаем близкими?}
\textbf{Время:} 01:15

\textbf{Тезис:} пространственная разность q описывает относительное движение, а полусумма K — общее.

\textbf{Визуал:} два вектора импульса; q соединяет их концы, K направлен между ними.

\textbf{Текст:} Пока достаточно обычных трёхмерных векторов. Вычтем второй импульс из первого и назовём разность q. Близкие векторы дают малую разность. Не путайте близость импульсов с близостью направлений: одинакового угла недостаточно, если модули сильно отличаются.

Полусумму назовём K. Она описывает общее движение пары. Поперечный импульс k-T — модуль поперечной части этой векторной полусуммы. Это не среднее двух модулей: противоположные одинаковые векторы дают ноль.

Во всём рассказе q означает полную разность. Релятивистскую запись с энергией добавим тогда, когда она понадобится для фазы. Сейчас уже можно разбивать пары по относительному и среднему импульсам.

\textbf{Переход:} Теперь у нас есть переменные для пары. Объясним квантовую причину особой корреляции при малом q.

\clearpage
\section*{Слайд 09. Два назначения — один квантовый исход}
\textbf{Время:} 01:15

\textbf{Тезис:} неразличимые назначения импульсов тождественным бозонам складываются на уровне амплитуд.

\textbf{Визуал:} две точки и два импульса; пройти по прямым и обменным стрелкам.

\textbf{Текст:} Из точек A и B испущены два положительных пиона. Мы зарегистрировали импульсы p-один и p-два. Возможны два назначения: первый импульс относится к A или к B.

Для тождественных частиц перестановка меток не создаёт другого физически различимого конечного состояния. Это свойство квантового описания, а не просто недостаток нашего детектора.

Поэтому для бозонов складываем две амплитуды со знаком плюс. Амплитуды похожи на волны: результат зависит от относительной фазы. Вероятность получаем уже после сложения. Для положительного и отрицательного пионов именно этого обменного члена нет, хотя взаимодействия между ними тоже могут давать корреляции.

\textbf{Переход:} Запишем эти амплитуды и увидим, как разность импульсов связывается с разделением точек.

\clearpage
\section*{Слайд 10. Сначала амплитуды. Потом вероятность}
\textbf{Время:} 02:30

\textbf{Тезис:} обменная симметризация даёт вес 1 + cos фазы, содержащей разность импульсов и разделение точек.

\textbf{Визуал:} правило вероятности → амплитуды → раскрытый квадрат → фаза.

\textbf{Текст:} Амплитуду можно представить комплексной стрелкой: длина — модуль, угол — фаза. Вероятность пропорциональна квадрату модуля. Звёздочка означает комплексное сопряжение, Re — действительную часть. Этого достаточно для нашего вывода.

Пока используем свободные плоские волны после испускания и одинаковые модули амплитуд для близких импульсов. Прямое и обменное назначения дают a и b. Складываем их и делим на корень из двух.

Раскрываем квадрат. Диагональные части после деления дают единичную базовую линию. Два перекрёстных слагаемых сопряжены; их сумма — удвоенная действительная часть произведения a на b-звёздочка. Получается единица плюс косинус разности фаз.

Теперь добавим энергии и времена. Ранее q был пространственной разностью. Четырёхимпульс объединяет энергию и обычный импульс, а его полная разность содержит ещё q-ноль — разность энергий. Четырёхкоордината объединяет ct и положение. Фаза содержит q-ноль на разность ct минус пространственное скалярное произведение. Разделив на h-bar-c, получаем безразмерный угол при импульсах в ГэВ и расстояниях в фемтометрах.

Пауза: пройти группировку четырёх исходных фаз в одно произведение. При нулевом q вес равен двум.

\textbf{Переход:} Мы вывели вес для двух точек. Что получится после усреднения по источнику?

\clearpage
\section*{Слайд 11. От двух точек — ко всему источнику}
\textbf{Время:} 01:45

\textbf{Тезис:} независимое усреднение двух точек приводит к квадрату модуля Фурье-образа источника.

\textbf{Визуал:} повторять выбор A и B → двойное среднее → факторизация.

\textbf{Текст:} Плотность rho описывает вероятности точек и времён испускания при выбранном среднем импульсе K. Её интеграл равен единице.

Усредняем косинус по всем парам точек. При независимом испускании вероятность пары — произведение двух плотностей, поэтому возникает двойной интеграл. Экспонента разности фаз раскладывается на множитель для A и множитель для B. Интегралы оказываются сопряжёнными; произведение равно квадрату модуля одного интеграла. Так появляется преобразование Фурье.

Нужны хаотическое испускание, плоские волны и гладкая импульсная зависимость источника, позволяющая использовать общий K. При малом q фазы мало различаются; при большом взаимно усредняются. Большим расстояниям достаточно меньшего изменения q для такой компенсации.

Пауза: модуль теряет фазу, а энергия связана с импульсом. Поэтому произвольное четырёхмерное изображение источника однозначно не восстанавливается.

\textbf{Переход:} Проверим смысл этой связи на источнике с простой известной формой.

\clearpage
\section*{Слайд 12. Больше источник — уже пик}
\textbf{Время:} 02:00

\textbf{Тезис:} характерная импульсная ширина обратно пропорциональна размеру гауссова источника.

\textbf{Визуал:} две учебные кривые одной высоты; сравнить именно их ширину.

\textbf{Текст:} Возьмём мгновенный изотропный гауссов источник. R — ширина одной декартовой координаты. Это не среднеквадратичное расстояние от центра: там появился бы корень из трёх.

После преобразования Фурье и квадрата модуля получаем гауссову корреляцию. Её ширина имеет масштаб h-bar-c, делённый на R. Здесь размеры три и шесть фемтометров заданы заранее. Для них характерные ширины примерно 66 и 33 МэВ на c. Удвоили источник — вдвое сузили импульсный сигнал.

Это учебные кривые, не данные STAR и не аппроксимация SMASH. Высоту задаёт lambda. На разрешённую силу сигнала влияют, например, долгоживущие распады, примеси и разрешение; при обсуждении когерентности нужны дополнительные предположения. По одной высоте размер читать нельзя.

Пауза: попросить выбрать более крупный источник, не подглядывая в легенду.

\textbf{Переход:} Связь C с размером получена. Теперь разберёмся, как измерить импульсы, выбрать пионы и построить эту функцию из настоящих сигналов детектора.

\clearpage
\section*{Слайд 13. TPC: как частица оставляет трёхмерный след}
\textbf{Время:} 01:45

\textbf{Тезис:} время дрейфа электронов восстанавливает координату трека, а не скорость пролетающей частицы.

\textbf{Визуал:} частица → ионизационные электроны → дрейф к торцу → считывающие площадки.

\textbf{Текст:} Чтобы измерить обещанную C, сначала нужны надёжно восстановленные импульсы. Начнём с трекинга в газовой камере. Заряженная частица ионизует атомы: вдоль её пути остаются электроны. По ним нужно восстановить пространственный след.

Электрическое поле переносит электроны к считывающим плоскостям на торцах камеры. В двух половинах TPC дрейф направлен к противоположным торцам. Положение сработавших площадок даёт две координаты, а время прихода заряда — третью. Если скорость дрейфа известна, произведение этой скорости на время дрейфа даёт расстояние от места ионизации до плоскости считывания.

Здесь важно различить две скорости. Пион уже пролетел через газ; мы ждём созданные им электроны. Время их дрейфа говорит о координате следа, а не о скорости пиона. Собрав множество таких пространственных точек и связав их в последовательность, восстанавливаем трек.

Пауза: пройти указкой один участок — от места ионизации до сигнала на торце.

\textbf{Переход:} По треку теперь можно определить импульс. Как именно это делается?

\clearpage
\section*{Слайд 14. По кривизне — импульс, по ионизации — подсказка}
\textbf{Время:} 01:45

\textbf{Тезис:} TPC даёт импульс и информацию для идентификации, но одного признака не всегда достаточно.

\textbf{Визуал:} сначала траектории в магнитном поле, затем реальные полосы ионизационных потерь.

\textbf{Текст:} В магнитном поле заряженная частица отклоняется. При одном поле более жёсткая траектория означает больший поперечный импульс на единицу заряда. В практической формуле поле берётся в теслах, радиус кривизны — в метрах, а результат получается в ГэВ на c. Q/e — заряд в единицах элементарного; для наших пионов его модуль равен единице. Направление изгиба помогает определить знак заряда.

Кривизна в поперечной плоскости даёт p-T. Наклон пространственного трека добавляет продольную компоненту, поэтому восстанавливается полный вектор импульса.

Но импульс ещё не говорит, какая частица прошла. Дополнительная подсказка — средние потери энергии на ионизацию вдоль трека. На настоящем графике разные сорта частиц образуют полосы: π — пионы, K — каоны, p — протоны. При одинаковом импульсе разные массы означают разные скорости и разные ожидаемые потери. Полосы имеют ширину и могут перекрываться; нельзя считать каждый цвет безошибочной меткой.

\textbf{Переход:} Импульс известен; добавим время пролёта и уточним, какая частица зарегистрирована.

\clearpage
\section*{Слайд 15. TOF: измеряем скорость и узнаём массу}
\textbf{Время:} 01:45

\textbf{Тезис:} сочетание траектории, времени пролёта и импульса позволяет идентифицировать частицу.

\textbf{Визуал:} отметить начало отсчёта, путь L, приход в TOF и разделение полос по массе.

\textbf{Текст:} TOF означает время пролёта. Из восстановленной траектории берём длину пути L до детектора. Вычитаем момент начала t-ноль из времени прихода сигнала и получаем время движения частицы. Оценка t-ноль нужна отдельно: одно время срабатывания TOF без начала отсчёта ещё не даёт скорость.

Делим длину пути на время и на скорость света — получаем beta. Теперь используем импульс из TPC. Из релятивистской связи энергии, импульса и массы возникает формула для массы в квадрате; в этой записи скорость света положена равной единице. Специальный вывод пока не нужен: два дополняющих друг друга измерения, импульс и скорость, позволяют различать массы.

На графике видны полосы, а не точные линии: приборы имеют разрешение. В анализе выбирают пионную область с учётом TPC и TOF. Итак, нам понадобились оба прибора: один восстановил движение в пространстве, второй добавил время пролёта и помог узнать сорт частицы.

\textbf{Переход:} Теперь у нас есть события с пионными кандидатами и их импульсами. Как превратить этот список в измеренную C?

\clearpage
\section*{Слайд 16. Пары и контроль: как построить корреляцию}
\textbf{Время:} 02:00

\textbf{Тезис:} отношение пар одного события к смешанному контролю оценивает корреляционный избыток.

\textbf{Визуал:} распределения A и B → нормированное отношение; пояснение центральности дать при слове «похожие».

\textbf{Текст:} В A помещаем пары пионов одного события. Его форма зависит и от корреляций, и от спектров и регистрации. Для контроля берём частицы из разных похожих событий и получаем B.

Почему похожих? Центральность — шкала, связанная с геометрией перекрытия ядер. В эксперименте её оценивают, например, по множественности заряженных частиц и калибруют по долям полного адронного сечения. Класс 0–10 процентов — наиболее центральная десятая часть этого сечения, а не десять процентов радиуса ядра. Внутри близких классов смешиваем события с сопоставимыми спектрами. Учитываем также вершину и условия регистрации.

Нормированное отношение A к B оценивает введённое ранее сравнение плотностей выходов частиц и пар. Общая нормировка задаёт C = 1 для независимого фона в контрольной области, где фемтоскопический эффект мал. Все парные эффекты и нефемтоскопические связи смешивание автоматически не устраняет.

На учебном рисунке два счёта сгенерированы независимо по Пуассону, а событий и настоящего смешивания нет. Экспозиция B больше в десять раз; нормировка десять задана заранее. Пример объясняет деление, но не зависимости реального фона.

\textbf{Переход:} Теперь проверим, насколько точно мы знаем полученное отношение.

\clearpage
\section*{Слайд 17. 4950 пар — не 4950 независимых измерений}
\textbf{Время:} 01:30

\textbf{Тезис:} статистическая точность зависит от структуры событий, а не только от числа пар.

\textbf{Визуал:} один общий трек выделен сразу в нескольких парах; затем цепочка повторного анализа.

\textbf{Текст:} Сто треков дают 4950 неупорядоченных пар. Но один трек используется много раз. Если с ним что-то изменится, изменятся сразу несколько пар. Поэтому нельзя считать все записи независимыми измерениями.

Пуассоновская оценка через квадратный корень из счёта полезна как первое приближение. Формула ошибки отношения требует независимых A и B и фиксированной нормировки. Общие треки, повторное смешивание и оценённая по данным нормировка создают зависимости.

Для проверки можно пересобирать выборку целыми событиями и повторять анализ: фон, нормировку и аппроксимацию. Это bootstrap по событиям. Отдельно меняют отборы и диапазоны анализа, проверяя систематические смещения. Раздвоение трека создаёт ложную пару, а слияние близких треков теряет настоящую. Большая статистика не исправляет такую проблему автоматически.

\textbf{Переход:} Прежде чем показывать данные STAR, определим три направления, по которым исследуется ширина корреляции.

\clearpage
\section*{Слайд 18. Три направления вместо одного радиуса}
\textbf{Время:} 01:30

\textbf{Тезис:} направленные радиусы отделяют продольную и поперечную структуру области испускания.

\textbf{Визуал:} сначала пучок и long, затем K-T и out, затем side.

\textbf{Текст:} Long направлено вдоль пучка. Out — вдоль среднего поперечного импульса конкретной пары. Side перпендикулярно обоим. Это взаимно перпендикулярные пространственные направления, даже если на рисунке они показаны в перспективе.

Для анализа переходим в LCMS: убираем суммарное продольное движение пары преобразованием Лоренца. Поперечное движение сохраняется, поэтому LCMS не является системой покоя пары.

Теперь вместо одной ширины аппроксимируем структуру по трём компонентам q и получаем три радиуса. Диагональная гауссова форма — приближение с определёнными симметриями. В более общем случае нужны смешанные члены; пример вынесен в резерв.

Мы определили язык результата. Пора посмотреть, как выглядит настоящая измеренная корреляция.

\textbf{Переход:} Посмотрим на настоящую трёхмерную корреляцию и на взаимодействия, которые меняют её форму.

\clearpage
\section*{Слайд 19. Так выглядит корреляция STAR при 3 ГэВ}
\textbf{Время:} 01:30

\textbf{Тезис:} экспериментальная форма требует учёта кулоновского взаимодействия пары.

\textbf{Визуал:} реальные проекции рисунка 4 статьи Краевой; затем соответствующий кулоновский множитель.

\textbf{Текст:} Это измерение Анны Краевой от имени STAR: центральность 0–10 процентов, k-T от 0,15 до 0,25 ГэВ на c. Цвета различают положительные и отрицательные пионы. В каждой проекции остальные компоненты q ограничены интервалом плюс-минус 0,05 ГэВ на c. Перед нами проекции одной трёхмерной функции.

У одинаковых зарядов есть отталкивание. Оно меняет корреляцию при малом относительном импульсе и конкурирует с обменным усилением. Для аппроксимации применена процедура Боулера–Синюкова с кулоновским фактором. Поэтому свободная гауссова формула сама по себе не описывает эти точки.

В других каналах существенно и сильное взаимодействие. Общий принцип — использовать волновую функцию пары с нужными взаимодействиями и усреднять её по источнику.

\textbf{Переход:} Чтобы интерпретировать извлечённые радиусы, нужно уточнить: какую часть расширяющегося источника выбирает данная пара?

\clearpage
\section*{Слайд 20. Почему быстрые пары видят меньшую область?}
\textbf{Время:} 01:30

\textbf{Тезис:} поток связывает координаты и импульсы, выделяя области однородности.

\textbf{Визуал:} общий источник → стрелки расширения → выделенный участок для выбранного K.

\textbf{Текст:} Расширяющееся вещество в среднем движется наружу. Координата испускания и импульс оказываются связаны, хотя тепловое движение эту связь размывает.

Отберём пары, летящие вправо с определённым средним импульсом. Они чаще приходят оттуда, где движение среды согласовано с их импульсом. Другому направлению соответствует другой участок. Это область однородности — часть всего источника, выделенная нашим отбором.

При увеличении k-T такая область может сужаться. Поэтому радиус уменьшается, хотя общий объём не обязан уменьшаться. Схема объясняет возможный механизм, а не определяет единственную динамическую модель.

Теперь понятно, зачем мы всё время сохраняли K и классы пар: без них само слово «размер» недостаточно определено.

\textbf{Переход:} С этой оговоркой посмотрим, как извлечённые радиусы STAR меняются со средним импульсом пары.

\clearpage
\section*{Слайд 21. Что STAR узнаёт из зависимости радиусов?}
\textbf{Время:} 01:45

\textbf{Тезис:} зависимости радиусов проверяют пространственную структуру и динамику расширения выбранного источника.

\textbf{Визуал:} направления и заряды в легенде → ошибки → изменение с k-T.

\textbf{Текст:} Рисунок показывает три LCMS-радиуса STAR для наиболее центральных столкновений. Цвет отвечает за направление, заполнение маркера — за заряд пионов. Сначала прочитаем эти обозначения, чтобы не сравнивать разные направления как одну последовательность.

По мере роста k-T радиусы уменьшаются. Это согласуется с картиной потока и областей однородности: меняется часть источника, из которой преимущественно приходят выбранные пары. Такое измерение ограничивает возможные сценарии расширения, но не выбирает единственную модель без дополнительных наблюдаемых.

Кроме статистических ошибок показана систематическая неопределённость. В этом варианте анализа не было отдельного отбора по быстроте пары; доступная выборка всё равно ограничена регистрацией.

Различия зарядов могут отражать несколько механизмов, включая распады и электрическое поле среды. Для их объяснения одной геометрической фразы недостаточно.

\textbf{Переход:} Размеры чувствительны и к динамике. Можно ли дополнительно прочитать по ним длительность испускания?

\clearpage
\section*{Слайд 22. Почему радиус out не является прямыми часами?}
\textbf{Время:} 01:30

\textbf{Тезис:} временная дисперсия смешивается с пространственной и с их ковариацией.

\textbf{Визуал:} три слагаемых в R-out; особо выделить среднее.

\textbf{Текст:} Для одинаковых масс энергия и импульс связаны так, что q на K равно нулю. Поэтому в интерференционной фазе возникает комбинация положения и времени: x-out минус скорость пары на ct.

Разброс этой комбинации даёт радиус out. В раскрытом квадрате видны три части: пространственная дисперсия, ковариация координаты со временем и временная дисперсия с коэффициентом скорости в квадрате.

Средний член отвечает на вопрос: связано ли более позднее испускание с другим положением? Он может компенсировать или усиливать влияние времени. Поэтому отношение out к side нельзя прочитать как отдельные часы.

Пауза: чтобы разность квадратов двух радиусов дала только временной вклад, нужны дополнительные условия на геометрию и ковариацию. Их проверка — задача интерпретации.

\textbf{Переход:} Модель с известной историей помогает исследовать именно такие условия.

\clearpage
\section*{Слайд 23. SMASH: проверяем метод на известной истории}
\textbf{Время:} 01:15

\textbf{Тезис:} транспортный расчёт даёт возможность сопоставить источник и результат его корреляционного анализа.

\textbf{Визуал:} заданное столкновение → отдельные адронные взаимодействия → история частиц.

\textbf{Текст:} SMASH моделирует движение и реакции адронов. Частицы распространяются, сталкиваются, образуют резонансы и распадаются. Получается набор событий с импульсами, а при соответствующем выводе истории — с координатами и временами взаимодействий.

В эксперименте пространственная история неизвестна, и мы оцениваем её характеристики по корреляциям. В модели можно поставить проверку: взять известный источник, обработать его тем же методом и выяснить, что именно восстановилось.

Это адронный транспорт. В гибридном описании он может отвечать за позднюю адронную стадию после гидродинамики. Применимость конкретного расчёта проверяется данными и настройками модели.

\textbf{Переход:} Однако список классических частиц ещё не создаёт выведенный нами обменный член.

\clearpage
\section*{Слайд 24. Как добавить квантовую интерференцию к событиям?}
\textbf{Время:} 02:00

\textbf{Тезис:} средний квантовый вес модельных пар отличается от экспериментального отношения к mixed-event фону.

\textbf{Визуал:} отобранная пара → фаза → вес → среднее в импульсном бине.

\textbf{Текст:} Для каждой пары берём импульсы и точки испускания. Из их разностей получаем фазу и назначаем вес: единица плюс косинус. Складываем веса в бине и делим на число тех же пар. Результат — средний вес.

Здесь нет двух независимых счётчиков. В числитель и знаменатель входят одни и те же пары, поэтому это другая конструкция, чем экспериментальное смешивание событий. Ошибку отношения независимых пуассоновских гистограмм использовать нельзя.

Нужны точки последнего взаимодействия или рождения конечного пиона при распаде. Общая поздняя поверхность вывода содержит уже свободный пролёт; её координаты не равны точкам испускания. Формат данных генератора необходимо проверить.

Косинус учитывает свободные плоские волны. Для кулоновского и сильного взаимодействий нужен другой вес с соответствующей волновой функцией. Начать проверку полезно с искусственного гауссова источника известной ширины, затем переходить к транспортным событиям. Неопределённость итоговых параметров проверяют с учётом структуры событий.

\textbf{Переход:} Посмотрим на конкретный результат моего освоения этой обработки.

\clearpage
\section*{Слайд 25. Мой расчёт: выбранные пары → корреляция}
\textbf{Время:} 01:45

\textbf{Тезис:} модельная корреляция меняется вместе с определением класса пар.

\textbf{Визуал:} крупная исходная корреляция из защиты 2025 года; справа условия и интервалы k-T.

\textbf{Текст:} Это оригинальный график из моей защиты 2025 года, а не новый расчёт к сегодняшней лекции. Анализировались модельные пионы SMASH в Au+Au при трёх ГэВ. На рисунке один класс центральности, 0–10 процентов, и четыре интервала k-T.

По горизонтали q-inv — инвариантная мера различия импульсов. В системе покоя пары она равна удвоенному модулю импульса одного пиона. Она не совпадает с модулем пространственной разности в произвольной системе; полное определение оставлено в резерве. По вертикали — корреляционная функция.

При большом q она приближается к единице, при малом виден сигнал. Меняя средний импульс пары, мы меняем усреднение по источнику. Поэтому разные кривые имеют физический смысл лишь вместе с указанными отборами.

Точки и ошибки сохранены из архива. Полного описания неопределённостей там нет, поэтому не будем оценивать по картинке значимость каждого различия. Здесь важно увидеть всю цепочку обработки от событий к корреляции.

\textbf{Переход:} Следующий шаг — выразить форму через параметры аппроксимации.

\clearpage
\section*{Слайд 26. Что получилось из аппроксимации?}
\textbf{Время:} 02:00

\textbf{Тезис:} архивный R-inv — эффективный параметр одномерного анализа с установленными ограничениями проверки.

\textbf{Визуал:} большой оригинальный график R-inv; показать центральности и исходные ошибки.

\textbf{Текст:} По горизонтали растёт k-T, по вертикали — R-inv в фемтометрах. В более центральных классах значения в целом выше. У самого центрального класса виден убывающий тренд, но он не обязан одинаково проявляться в каждом классе и каждом следующем бине.

Этот параметр получен одномерной аппроксимацией по q-inv. Он не является одним из трёх LCMS-радиусов и не обязан равняться их среднему. Усреднение по направлениям и движение пары меняют связь с источником.

Есть существенная граница достоверности. В формулах архивных слайдов встречаются ошибки записи, а код расчёта недоступен. Метод вычисления исходных ошибок также не установлен. Поэтому численные радиусы в выбранной конвенции независимо не подтверждены. Показываем сохранённый результат выполненной тогда работы, а не новую верификацию его точности.

Пауза: для количественного физического вывода нужно вернуться к событиям, проверить аппроксимацию и неопределённости.

\textbf{Переход:} Что потребуется, чтобы такой модельный результат действительно проверить экспериментом?

\clearpage
\section*{Слайд 27. Как проверить, описывает ли модель эксперимент?}
\textbf{Время:} 01:00

\textbf{Тезис:} сравнение должно относиться к одинаково определённым наблюдаемым и классам пар.

\textbf{Визуал:} сопоставить строки таблицы эксперимента и модели.

\textbf{Текст:} Одинаковых Au+Au и трёх ГэВ недостаточно. Нужно согласовать заряды, интервалы импульса, центральность и продольные отборы. Быстрота пары и псевдобыстрота отдельной частицы — разные ограничения.

Далее выбираем одинаковую систему отсчёта и размерность анализа. Одномерный R-inv нельзя напрямую считать трёхмерным радиусом STAR. Проверяем взаимодействия: у показанного эксперимента есть учёт Кулона, простой косинус модели его не включает.

Наконец, нужны проверенные ошибки, нормировка и отбор событий. Тогда различие сможет ограничить физическую модель, а не просто показать, что две программы посчитали разные величины.

\textbf{Переход:} Вернёмся к исходному вопросу о размерах источника, который уже исчез.

\clearpage
\section*{Слайд 28. Детектор → корреляция → источник: что и зачем?}
\textbf{Время:} 01:15

\textbf{Тезис:} свойства источника исследуются через измеримые корреляции и проверяемую модель их связи с пространством и временем.

\textbf{Визуал:} пройти всю цепочку и связать каждый переход с его вопросом.

\textbf{Текст:} Детекторы дали треки, импульсы и идентификацию. Статистика пар добавила сведения, которых нет в отдельных спектрах. Бозонная интерференция связала относительный импульс с разделением точек испускания.

Так мы получили доступ к пространственно-временным масштабам выбранной области источника. Они помогают проверять сценарии расширения и работу транспортных моделей. Для вывода о динамике нужны те же пары, согласованные отборы и правильное описание взаимодействий.

Оставим короткий вопрос аудитории: в простом гауссовом примере размер увеличили вдвое. Как изменилась ширина импульсного сигнала? Дать ответить: уменьшилась вдвое. Затем предложить обсудить, какие дополнительные предположения потребовались бы для перехода от такого размера к длительности.

\clearpage
\section*{Слайд 29. Источники: теория, данные STAR и транспорт}
\textbf{Время:} резерв

\textbf{Тезис:} формализм, измерение и модель имеют самостоятельные первоисточники.

\textbf{Визуал:} теория → статья STAR → исходная работа и документация SMASH.

\textbf{Текст:} Основной теоретический обзор — Lisa, Pratt, Soltz и Wiedemann, 2005 год, arXiv:nucl-ex/0505014. Формулы лекции записаны в конвенции полной разности q, поэтому соглашения обзора пересчитаны.

Измерения при трёх ГэВ взяты из статьи Анны Краевой от имени STAR: Universe 2024, том 10, статья 188, DOI 10.3390/universe10040188. У доклада и опубликованной статьи отдельные даты и статусы; исторические пометки STAR Preliminary относятся к конкретному представленному материалу.

SMASH описана в работе Weil и соавторов 2016 года и официальной документации. Приведённые модельные события автора относятся к архивной работе 2025 года; номер современного релиза генератора им не приписывается.

\clearpage
\section*{Слайд 30. Откуда взялась идея и как работают детекторы}
\textbf{Время:} резерв

\textbf{Тезис:} современный метод соединяет общую квантовую основу и развитие способов пространственно-временного анализа.

\textbf{Визуал:} исторические первоисточники и отдельные работы об измерении треков и времени пролёта.

\textbf{Текст:} В 1956 году Хэнбери Браун и Твисс опубликовали лабораторную демонстрацию корреляций фототоков: Nature, том 177, страница 27, DOI 10.1038/177027a0. Работа о Сириусе опубликована в Nature, том 178, страница 1046, DOI 10.1038/1781046a0. Без полного разбора исходных данных не называем точные параметры установки или диаметр звезды.

Пионные экспериментальные результаты обсуждались уже в 1959 году. Поэтому работу Goldhaber, Goldhaber, Lee и Pais 1960 года корректнее связывать с объяснением роли бозонной симметризации в угловых корреляциях при аннигиляции антипротона с протоном, а не называть первым измерением тяжёлоионного радиуса. Размер модельного объёма уже входил в этот расчёт; связь с размером не появилась впервые только в 1970-е годы.

Работы Копылова и Подгорецкого 1970-х годов развивали связь корреляций с координатами и временами испускания. Оптическая база приёмников и относительный импульс пионов — разные экспериментальные переменные; общий принцип не делает установки одинаковыми.

Детекторные схемы опираются на работу Anderson и соавторов о TPC STAR, а совместная идентификация — на статьи Shao и соавторов и Llope и соавторов. При вопросе о времени вернуться к различию дрейфа электронов, пролёта частицы и отдельной оценки начала события.

\clearpage
\section*{Слайд 31. Авторские источники: условия и границы проверки}
\textbf{Время:} резерв

\textbf{Тезис:} архив подтверждает происхождение изображений и перечисленные условия, но не полный воспроизводимый расчёт.

\textbf{Визуал:} таблица отборов и указание страниц защиты.

\textbf{Текст:} В 2023 году подготовлена обзорная курсовая; собственные численные результаты моделирования ей не приписываются. Полный текст доступен. Ссылка за 2025 год открывает презентацию защиты на 28 страницах.

Для расчёта 2025 года указаны Au+Au при трёх ГэВ; пары одинаково заряженных пионов; полный и поперечный импульсы от 0,15 до 1,5 ГэВ на c; псевдобыстрота от минус 1,5 до 1,5. Интервалы k-T: 0,15–0,25; 0,25–0,35; 0,35–0,45; 0,45–0,60 ГэВ на c. Центральности: 0–10, 10–30, 30–50 и 50–80 процентов.

Заявлено исключение упругих событий, но программный критерий не установлен. Неизвестные число событий, версия SMASH, настройки аппроксимации и метод ошибок не дополнены данными из отдельного отчёта 2026 года. Ошибки записи архивных формул требуют проверки исходной реализации.

\clearpage
\section*{Слайд 32. Средний вес имеет другую статистическую ошибку}
\textbf{Время:} резерв

\textbf{Тезис:} общие пары связывают числитель и знаменатель и требуют учёта структуры событий.

\textbf{Визуал:} среднее независимых весов → ковариация отношения → пересборка событий.

\textbf{Текст:} Если сумма весов делится на число тех же пар, это выборочное среднее. При независимых весах и фиксированном n его дисперсия оценивается выборочной дисперсией весов, делённой на n. В раскрытой записи возникает знаменатель n на n минус один.

Если все веса единичные, средний вес в точности единица. Условная ошибка по весам нулевая — в отличие от результата механической подстановки в формулу двух независимых пуассоновских счётов.

Для общего отношения нужен ковариационный член. Однако его наличие ещё не учитывает повторение треков в разных парах. Конечные неопределённости проверяют на уровне событий с повторением полного анализа. Оценённая из данных нормировка тоже приносит дисперсию и ковариации.

\clearpage
\section*{Слайд 33. Архив 2025: трёхмерные радиусы out и side}
\textbf{Время:} резерв

\textbf{Тезис:} направленные параметры сохраняют информацию, потерянную при одномерном усреднении.

\textbf{Визуал:} две увеличенные оригинальные панели из слайда 17 защиты.

\textbf{Текст:} Эти графики подтверждают, что в работе выполнялась и трёхмерная аппроксимация. Точки и исходные ошибки не изменены, нового расчёта при подготовке лекции нет.

Out и side отвечают разным комбинациям координат и времени. Поэтому их нельзя заменить одним средним размером без дополнительных предположений. Для сравнения со STAR нужно подтвердить систему отсчёта, отборы, параметризацию, необходимые перекрёстные члены и способ расчёта ошибок.

Архивный график показывает проделанный этап анализа. Он не снимает вопросов к воспроизводимости и не заменяет проверку исходной реализации, обсуждавшуюся на основном слайде 26.

\clearpage
\section*{Слайд 34. Архив 2025: R-long и сила 1D-корреляции}
\textbf{Время:} резерв

\textbf{Тезис:} продольный радиус и lambda относятся к конкретной процедуре анализа.

\textbf{Визуал:} отдельно назвать, какая панель взята из 3D-аппроксимации, а какая из 1D.

\textbf{Текст:} R-long характеризует продольную структуру выбранной области испускания в принятой системе отсчёта. На него влияют отбор и динамика; это не исходная длина ядра.

Lambda справа относится к одномерной аппроксимации. Этот параметр нельзя автоматически приравнивать силе сигнала другого фита. Для экспериментальной выборки важны примеси и разрешение. В чистом модельном наборе такие эффекты не возникают сами собой: нужно отдельно исследовать распады, форму источника и диапазон аппроксимации.

Ошибки оставлены исходными. Без событий и установленного алгоритма неопределённостей не будем давать новую численную оценку значимости архивных трендов.

\clearpage
\section*{Слайд 35. А если менять продольное движение пары?}
\textbf{Время:} резерв

\textbf{Тезис:} отбор по продольному движению проверяет симметрии и динамику источника.

\textbf{Визуал:} формула быстроты → оригинальная панель рисунка 6(e) статьи Краевой.

\textbf{Текст:} Быстрота пары строится из суммарных энергии и продольного импульса. Она задана в общей системе центра масс столкновения; затем для каждой пары отдельно строится LCMS. Это разные операции. Псевдобыстрота отдельного трека определяется углом и не заменяет быстроту пары.

В гауссовой экспоненте может сохраняться смешанный член out–long. Обозначение R-out–long в квадрате традиционно, но коэффициент имеет знак: это не квадрат положительной длины.

На графике коэффициент меняет знак около нулевой быстроты. Серые точки — отражение измеренных результатов, а не новые независимые данные; знак смешанного коэффициента при отражении тоже меняется. Авторы обсуждают признаки нарушения приближённой продольной буст-инвариантности. Один график не следует превращать в доказательство единственной модели эволюции.

\clearpage
\section*{Слайд 36. Полная конвенция импульсов: для вопросов после доклада}
\textbf{Время:} резерв

\textbf{Тезис:} инвариантный относительный импульс отличается от пространственной разности в произвольной системе.

\textbf{Визуал:} полная разность и полусумма; система покоя пары рядом с LCMS.

\textbf{Текст:} Четырёхимпульс содержит энергию и обычный импульс. Здесь q — полная разность, K — полусумма, метрика имеет знаки плюс, минус, минус, минус. В обзоре Lisa относительный импульс для одинаковых масс вдвое меньше нашего.

Q-inv определяется как корень из пространственной разности в квадрате минус разность энергий в квадрате. Для равных масс в системе покоя пары энергии одинаковы, и q-inv равен удвоенному модулю импульса одной частицы. В LCMS поперечное движение пары остаётся, поэтому системы не совпадают.

Средний поперечный импульс определяется через векторную полусумму. Смена конвенции q без пересчёта фазы и экспоненты меняет масштаб извлечённого радиуса.

\clearpage
\section*{Слайд 37. STAR при 3 ГэВ: неподвижная мишень и отбор частиц}
\textbf{Время:} резерв

\textbf{Тезис:} лабораторная геометрия и энергия в системе центра масс должны различаться явно.

\textbf{Визуал:} пучок золота → золотая мишень внутри STAR; положение области взаимодействия относительно TPC.

\textbf{Текст:} При фиксированной мишени один поток частиц налетает на покоящееся вещество. Система центра масс движется относительно лаборатории. Поэтому энергия столкновения пары нуклонов и энергия либо импульс пучка — разные величины.

В рассмотренной статье STAR использована золотая мишень внутри установки, смещённая от её центра. Это меняет доступную геометрию регистрации и продольное окно. Отбор общей вершины помогает выбрать нужные столкновения, но сам по себе не восстанавливает фемтометровые точки последнего взаимодействия пионов.

При моделировании нужно воспроизвести соответствующие преобразования и отборы. Одной подписи «Au+Au, 3 ГэВ» для сопоставления с данными недостаточно.

\clearpage
\section*{Слайд 38. Как STAR выбирает пионы и область наблюдения}
\textbf{Время:} резерв

\textbf{Тезис:} идентификация и окно регистрации являются частью определения экспериментального результата.

\textbf{Визуал:} реальные панели ионизационных потерь, массы и доступных пар из статьи Краевой.

\textbf{Текст:} Полосы ионизационных потерь в TPC сравнивают с ожидаемыми для разных частиц. TOF добавляет скорость и оценку массы. Совместные ограничения уменьшают примеси; в рассмотренном анализе оценённая чистота одиночной пионной выборки была не ниже 98 процентов. Это не утверждение об отсутствии всех детекторных эффектов.

Доступное окно пар неоднородно по среднему импульсу и продольному движению. Даже анализ без отдельного ограничения по быстроте не означает полный угловой охват. При сравнении с моделью нужно повторить реальные условия отбора и отдельно проверить эффекты близких треков.

Если на историческом рисунке точки зеркально отражены, это использование симметрии, а не расширение независимых измерений.

\end{document}
